This chapter presents the objectives of this thesis, the dataset used to verify whether these objectives are achieved, 
the experimental setup in which the tests are executed, and the tests performed, together with their results.

\section{Goals} \label{sec:goals}
    In order to prove that the instrumentation approach is valid, I have chosen three main goals:
    \begin{itemize}
        \item \textbf{Goal I} : The instrumented binary must be compatible with the original one.
            This means that the behavior of the program must be preserved after the instrumentation.
        \item \textbf{Goal II} : The instrumented binary must be compatible with AFL++ communication protocol.
        \item \textbf{Goal III} : Show that the instrumentation overhead is competitive with other binary-only rewriters 
                and lower than that of dynamic binary translators.
    \end{itemize}

    In order to prove that the tool satisfies these goals, the following experiments were conducted:
    \begin{itemize}
        \item \textbf{Experiment I} : Test the original binary and the instrumented binary with the same input, 
            without using a fuzzer, in order to prove that the behavior of the program is preserved after the instrumentation 
            (Goal I).
        \item \textbf{Experiment II} : Perform a fuzzing campaign with three different instrumentations:
        \begin{itemize}
            \item AFL++ in QEMU mode
            \item ZAFL 
            \item Mine
        \end{itemize}
        The chosen fuzzer is AFL++. This test aims to prove Goal II and Goal III.
    \end{itemize}

\section{Datasets and Experiment Setup}\label{sec:datasets_setup}
    \begin{center}
        \textbf{Experiment I}
    \end{center}
    The experiment consists in feeding the same input to both the original binary and the modified one, 
    and checking whether the outputs coincide. The dataset is taken from the \textit{real-world software} dataset used in the ZAFL 
    paper \cite{zafl}. Not all the software was used. \Cref{tab:experiment1Dataset} lists the binaries 
    used, with their versions.
    \begin{table}[ht]
        \centering
        \caption{Binaries extracted from ZAFL \textit{real-world software} dataset}
        \label{tab:experiment1Dataset}
        \begin{tabular}{c|c}
            Binary name & Version \\
            \hline
            bsdtar & bsdtar 3.7.2 - libarchive 3.7.2 \\
            \hline
            readelf & GNU Binutils 2.42\\
            \hline
            sfconvert & 0.3.6\\
            \hline
            tcpdump & 4.99.4\\
            \hline 
            unrtf & 0.21.10\\

        \end{tabular}
    \end{table}
    All the binaries considered for this experiment are stripped.
     \begin{center}
        \textbf{Experiment II}
    \end{center}
    The dataset used for the second experiment is LAVA-M \cite{lava}, in the same version used by ZAFL \cite{zafl}.
    This dataset is composed of four coreutils binaries (\textit{base64 , md5sum , uniq , who}), to which synthetic bugs 
    were added with the LAVA framework. Each bug has a unique ID and is protected by a checksum.
    An input can trigger a synthetic bug only if random mutation manages to satisfy the checksum.
    When a specific bug is triggered, the unique ID of that bug is printed and the program then crashes intentionally.\\ 
    The coreutils version used is \textbf{8.24}.
    As reported by Chen et al. \cite{Angora}, the LAVA-M version of \textit{md5sum} was modified incorrectly: 
    when run with the \texttt{-c} option, it crashes on every input. 
    This affects our second experiment, in which the fuzzing seed must be parsed with \texttt{-c}. 
    Since every input triggers a crash, the results for this binary are not meaningful.
    The experiment was run with AFL++ version 5.03c, over a 24-hour fuzzing campaign with three repetitions for each tested binary.
    The experiment was executed inside a KVM virtual machine (full virtualization) exposing 16 virtual CPUs, hosting 
    Ubuntu 24 as guest operating system, with 32 GB of RAM. 
    All the different runs was executed concurrently; each was pinned to a dedicated virtual CPU. 
    The 3 different instrumentation methods are tested indipendently; parallely are executed the 3 runs of one of the instrumentation 
    method on each of the 4 binaries, having a total of 12 virtual cores occupied.
    The seeds for each runs are different.
    The initial corpus is composed of only 1 test case.\\
    The general AFL++ command used to each fuzzing campaign is:\\
    \texttt{afl-fuzz -i <input\_dir> -o <findings\_dir> -V 86400   -\,- <binary>} ;
    with some difference depending on the binary that is being tested and the instrumentation method:
    \begin{itemize}
        \item For the binaries instrumented with ZAFL the environment variable \\``AFL\_SKIP\_BIN\_CHECK=1'' must be set
        \item For the binaries fuzzed with AFL++ in QEMU mode, the option ``-Q'' must be added to the command
    \end{itemize}
    Moreover the 4 binaries are executed in different ways:
    \begin{itemize}
        \item \textit{base64} : \texttt{base64 -d}
        \item \textit{md5sum} : \texttt{md5sum -c @@}
        \item \textit{uniq} : \texttt{uniq}
        \item \textit{who} : \texttt{who @@} 
    \end{itemize}

    
\section{Experiment I}\label{sec:experiment1}
Among the binaries listed in \cref{tab:experiment1Dataset}, 
\textit{tcpdump} and \textit{unrtf} could not be instrumented: the failure occurs 
during the GRAPE \gls{CFG} reconstruction (specifically during \texttt{Angr} CFG analysis).\\
Each binary is executed different times, with different inputs and different execution options.
The output of the original and the instrumented are then compared to check if are the same or not.
The results of the experiment are shown in \cref{tab:experiment1}.\\

    \begin{table}[ht]
        \centering
        \caption{Results of experiment I}
        \label{tab:experiment1}
        \begin{tabular}{c|c|c|c|c}
                        & Instrumented  & total tested inputs   & pass  & fail\\
            \hline
            bsdtar      & \cmark        &   34                  & 34    &   0\\
            \hline
            readelf     & \cmark        &   54                  & 4     &  50\\
            \hline
            sfconvert   & \cmark        &   52                  & 52    &   0\\
            \hline
            tcpdump     & \xmark        & \nd                   & \nd   & \nd\\
            \hline
            unrtf       & \xmark        & \nd                   & \nd   & \nd\\
        \end{tabular}
    \end{table}
    

    With regard to \textit{readelf}, the instrumented binary encounters a \texttt{SIGSEGV} error in all the failed tests. 
    This is due to the fact that, 
    when new sections and segments are added to the binary, the \textbf{lief} library automatically reorganizes the \gls{ELF} layout, 
    causing some sections (e.g. \textit{.bss}) to be shifted. The instructions that referenced these sections 
    will, after the layout change, reference wrong addresses, causing a \texttt{SIGSEGV} when they try to access them.\\
    The 4 test passed by readelf are 4 inputs that in the original binary causes the execution to stop with an exit code 
    different from 0. These inputs does not cause a \texttt{SIGSEGV} in the original as in the instrumented file, but 
    ``pass'' the test because in both case the exit code is different from 0. 
    A possible reason for the difficulty encountered in instrumenting this dataset could be the fact that this binaries are 
    stripped. The absence of symbols could create more problems in \gls{CFG} reconstruction.

\section{Experiment II}\label{sec:experiment2}
    Among the four LAVA-M binaries, none fails during the instrumentation phase.
    However, two of them (i.e. \textit{md5sum} and \textit{who}) result in an incorrect instrumentation:
    \textit{md5sum} was already corrupted, as stated before, while \textit{who} was incorrectly instrumented by the tool.
    The results of the experiment are shown in \cref{table:fuzzingStats}.
    The AFL++ in QEMU mode run and the ZAFL run of the \texttt{md5sum} and \texttt{who} are not shown in \cref{table:fuzzingStats} 
    because is impossible to compare those runs with the tool discussed in this thesis.
    \Cref{table:fuzzingStats} is therefore divided into two subtables, one per fuzzed binary.
    Each table reports some parameters obtained for each of the three runs performed.\\
    As can be seen, the tool developed in this thesis finds a considerably smaller number of edges than the other tested methods.
    This could be caused by a different edge instrumentation method (for example, because the tool does not instrument 
    function calls).\\
    Since each binary has been injected with LAVA synthetic bugs, the \textit{``uniq\_bugs''} contains the number of uniq bugs 
    triggered by the run.
    The id of each uniq bug are reported in \cref{tab:uniq_bugs_ids}.
    The number of uniq bugs triggered is always less or equal to the number of uniq crashes because different 
    crashes could trigger the same bug.
        
    ZAFL is the only tool capable of executing \textit{md5sum} with the \texttt{-c} option but the high number 
    of crashes if given by errors of the binary.
    Infact none of the found uniq crashes actually trigger a LAVA Bug, but trigger SIGSEGV errors that make the entire statistic 
    of md5sum unreliable.

    \Cref{tab:binariesStats} shows the size of the instrumented files in MBytes and the average executions per 
    second obtained from \cref{table:fuzzingStats}.
    As already discussed, the tool developed in this thesis produces larger binaries, because the \gls{SHM} update 
    is instrumented inline. As shown in \cref{tab:binariesStats} the size of all the instrumented binaries is similar 
    regardless of how many locations have been instrumented.
    This could be a sign of a fixed dimension incremention given by the added section.
    The average execution speed is only slightly higher than that of ZAFL.
    Since the tool execution speed, as can be seen in \cref{table:fuzzingStats}, is not always higher with respect to 
    the ZAFL execution speed, and only slightly higher in average, it is not possible to conclude an actual improvement in the execution speed.\\
    Moreover is present the speedup with respect to QEMU mode.
    The speedup is obtained dividing the average execution per second obtained with mine instrumentation and ZAFL by the 
    average executions per second obtanied with AFL++ in QEMU mode.

    Worth mention that ZAFL perform an optimization of the Intermediate Representation, that surely speed up the instrumented 
    binary execution. With inline instrumentation is possible to reach performance that are comparable to the one 
    obtained by ZAFL, without an actual optimization.
    \begin{table}[H]
        \centering
        \setlength{\tabcolsep}{3pt}
        \newcolumntype{L}{>{\raggedright\arraybackslash}p{2.6cm}}
        \newcolumntype{C}{>{\centering\arraybackslash}p{1.7cm}}
        \begin{subtable}{\linewidth}
            \centering
            \caption{\textmd{\texttt{base64}}}
            \label{table:fuzzingStats:base64}
            \resizebox{\linewidth}{!}{%
            \begin{tabular}{|L!{\vrule width 1.2pt}C|C|C!{\vrule width 1.2pt}C|C|C!{\vrule width 1.2pt}C|C|C|}
                \hline
                & \multicolumn{3}{c!{\vrule width 1.2pt}}{AFL++ QEMU mode} & \multicolumn{3}{c!{\vrule width 1.2pt}}{mine} & \multicolumn{3}{c|}{ZAFL} \\
                \hline
                & run1 & run2 & run3 & run1 & run2 & run3 & run1 & run2 & run3 \\
                \thickhline
                execs/s & 398 & 445 & 352 & 1706 & 1860 & 1776 & 1577 & 1792 & 1749 \\
                \hline
                edges\_found & 690 & 690 & 690 & 362 & 386 & 362 & 640 & 672 & 672 \\
                \hline
                corpus\_count & 162 & 169 & 178 & 151 & 214 & 149 & 143 & 205 & 201 \\
                \hline
                stability & 100\% & 100\% & 100\% & 100\% & 100\% & 100\% & 100\% & 100\% & 100\% \\
                \hline
                saved\_crashes & 7 & 3 & 2 & 5 & 102 & 13 & 10 & 115 & 118 \\
                \hline
                uniq\_bugs & 3 & 3 & 2 & 5 & 46 & 7 & 7 & 48 & 48 \\
                \hline
            \end{tabular}}
        \end{subtable}
        \bigskip
        \begin{subtable}{\linewidth}
            \centering
            \caption{\textmd{\texttt{uniq}}}
            \label{table:fuzzingStats:uniq}
            \resizebox{\linewidth}{!}{%
            \begin{tabular}{|L!{\vrule width 1.2pt}C|C|C!{\vrule width 1.2pt}C|C|C!{\vrule width 1.2pt}C|C|C|}
                \hline
                & \multicolumn{3}{c!{\vrule width 1.2pt}}{AFL++ QEMU mode} & \multicolumn{3}{c!{\vrule width 1.2pt}}{mine} & \multicolumn{3}{c|}{ZAFL} \\
                \hline
                & run1 & run2 & run3 & run1 & run2 & run3 & run1 & run2 & run3 \\
                \thickhline
                execs/s & 568 & 577 & 577 & 1751 & 1713 & 1783 & 1774 & 1560 & 1606 \\
                \hline
                edges\_found & 473 & 473 & 473 & 242 & 242 & 242 & 443 & 443 & 443 \\
                \hline
                corpus\_count & 97 & 101 & 100 & 86 & 88 & 87 & 90 & 89 & 85 \\
                \hline
                stability & 100\% & 100\% & 100\% & 100\% & 100\% & 100\% & 100\% & 100\% & 100\% \\
                \hline
                saved\_crashes & 2 & 0 & 0 & 0 & 0 & 0 & 0 & 1 & 0 \\
                \hline
                uniq\_bugs & 0 & 0 & 0 & 0 & 0 & 0 & 0 & 1 & 0 \\
                \hline
            \end{tabular}}
        \end{subtable}
        \bigskip
        
        \caption{Fuzzing statistics over a 24-hour fuzzing campaign on the LAVA-M dataset using AFL++ in QEMU mode, ZAFL instrumentation + AFL++ fuzzer, and my instrumentation + AFL++ fuzzer}
        \label{table:fuzzingStats}
        \vspace{2mm}
        \footnotesize
        The results for \texttt{md5sum} and \texttt{who} are not presented here because the comparison with the tool 
        discussed in this thesis is not possible    
    \end{table}
    \vspace{0.5cm}
    


    \begin{table}[htbp]
        \centering
        \begin{subtable}[t]{0.48\textwidth}
            \centering
            \caption{LAVA-M uniq bugs triggered by the runs in \textmd{\texttt{base64}}}
            \label{tab:uniq_base64}
            \begin{tabularx}{\linewidth}{|>{\raggedright\arraybackslash}p{1.5cm}|X|}
                \hline
                run name & uniq bugs IDs \\
                \hline
                afl1 & ['784', '792', '804']\\
                \hline
                afl2 & ['276', '792', '804']\\
                \hline
                afl3 & ['276', '790']\\
                \hline
                mine1 & ['784', '788', '792', '804', '815']\\
                \hline
                mine2 & [798', '804', '805', '806', '815', '817', '831', '835', '841', '842', '843']\\
                \hline
                mine3 & ['778', '782', '784', '788', '792', '804', '817']\\
                \hline
                zafl1 & ['253', '276', '560', '774', '778', '782', '792']\\
                \hline
                zafl2 & ['1', '222', '235', '253', '255', '274', '276', '278', '284', '386', '521', '526', '527', '554', '556', '558', '560', '562', '566', '572', '573', '576', '582', '583', '584', '774', '776', '778', '780', '782', '784', '786', '788', '790', '792', '798', '804', '805', '806', '813', '815', '817', '831', '832', '835', '841', '842', '843']\\
                \hline
                zafl3 & ['1', '222', '235', '253', '255', '274', '276', '278', '284', '386', '521', '526', '527', '554', '556', '558', '560', '562', '566', '572', '573', '576', '582', '583', '584', '774', '776', '778', '780', '782', '784', '786', '788', '790', '792', '798', '804', '805', '806', '813', '815', '817', '831', '832', '835', '841', '842', '843']\\
                \hline
            \end{tabularx}
        \end{subtable}
        \hfill
        \begin{subtable}[t]{0.48\textwidth}
            \centering
            \caption{LAVA-M uniq bugs triggered by the runs in \textmd{\texttt{uniq}}}
            \label{tab:uniq_uniq}
            \begin{tabularx}{\linewidth}{|>{\raggedright\arraybackslash}p{1.5cm}|X|}
                \hline
                run name & uniq bugs IDs \\
                \hline
                afl1 & []\\
                \hline
                afl2 & []\\
                \hline
                afl3 & []\\
                \hline
                mine1 & []\\
                \hline
                mine2 & []\\
                \hline
                mine3 & []\\
                \hline
                zafl1 & []\\
                \hline
                zafl2 & [472]\\
                \hline
                zafl3 & []\\
                \hline
            \end{tabularx}
        \end{subtable}

        \caption{LAVA-M uniq bugs triggered by the runs performed in \cref{table:fuzzingStats}}
        \label{tab:uniq_bugs_ids}
    \end{table}

    \begin{table}[ht]
        \centering
        \begin{tabular}{|c|c!{\vrule width 1.2pt}c|c!{\vrule width 1.2pt}c|}
            \hline
            &  & file\_size[MByte] & avg execs/sec & speedup w.r.t. AFL++ in QEMU mode\\
            \thickhline
            \multirow{2}{*}{base64}
            & mine & 8,32  & 1781 & 4,47\\
            \cline{2-5}
            & zafl & 0,14 & 1706  & 4,28\\
            \thickhline
            \multirow{2}{*}{md5sum}
            & mine & 8,32 & \xmark & \nd \\
            \cline{2-5}
            & zafl & 0,15 & 515 & \nd \\
            \thickhline
            \multirow{2}{*}{uniq}
            & mine & 8,32 & 1749 & 3,04\\
            \cline{2-5}
            & zafl & 0,15 & 1647 & 2,87\\
            \thickhline
            \multirow{2}{*}{who}
            & mine & 10,78  & \xmark & \nd \\
            \cline{2-5}
            & zafl & 1,72 & 456 & \nd \\
            \hline
        \end{tabular}
        \caption{Instrumented binaries statistics}
        \label{tab:binariesStats}
        \footnotesize
        \textit{Legend:}
        The ``\nd'' symbol is placed on that binaries where the comparison can not be done
    \end{table}